\section{第 3 章：无类型算术表达式}

\begin{taplauthorsolutionentrybox}
\phantomsection
\label{sol:author:3.2.4}
\textbf{3.2.4 解答}

\[
  \lvert S_{i+1}\rvert
  =\lvert S_i\rvert^3+3\lvert S_i\rvert+3,
  \qquad
  \lvert S_0\rvert=0
\]
因此 \(\lvert S_3\rvert=59439\)。

\taplsolutionbacklink{ex:3.2.4}{返回练习 3.2.4}
\end{taplauthorsolutionentrybox}

\begin{taplauthorsolutionentrybox}
\phantomsection
\label{sol:author:3.2.5}
\textbf{3.2.5 解答}

直接作归纳即可。\(i=0\) 时没有什么需要证明。以下设
\(i=j+1\)，其中 \(j\geq0\)，并以 \(S_j\subseteq S_i\) 为归纳假设；
我们要证明 \(S_i\subseteq S_{i+1}\)。任取 \(t\in S_i\)。由 \(S_i\)
是三个集合之并的定义，\(t\) 必为以下三种形式之一：

\begin{enumerate}
  \item \(t\in\{\tapltrue,\taplfalse,\taplzero\}\)。由 \(S_{i+1}\)
    的定义，显然 \(t\in S_{i+1}\)。
  \item \(t\) 为 \(\taplsucc{t_1}\)、\(\taplpred{t_1}\) 或
    \(\tapliszero{t_1}\)，其中 \(t_1\in S_j\)。由归纳假设，
    \(t_1\in S_i\)，故按定义 \(t\in S_{i+1}\)。
  \item \(t=\taplif{t_1}{t_2}{t_3}\)，其中
    \(t_1,t_2,t_3\in S_j\)。再次由归纳假设，三个子项都属于
    \(S_i\)，故按定义 \(t\in S_{i+1}\)。
\end{enumerate}

\taplsolutionbacklink{ex:3.2.5}{返回练习 3.2.5}
\end{taplauthorsolutionentrybox}

\begin{taplauthorsolutionentrybox}
\phantomsection
\label{sol:author:3.3.4}
\textbf{3.3.4 解答}

这里只证明对深度归纳的原理；另外两条与之类似。已知对每个项 \(s\)，只要
对所有满足 \(\termdepth(r)<\termdepth(s)\) 的项 \(r\) 都有 \(P(r)\)，
就能证明 \(P(s)\)；需要推出 \(P(s)\) 对所有 \(s\) 成立。定义自然数上的
新谓词
\[
  Q(n)\quad\Longleftrightarrow\quad
  \text{对每个满足 \(\termdepth(s)=n\) 的项 \(s\)，都有 \(P(s)\)。}
\]
现在对自然数 \(n\) 使用完全归纳原理（2.4.2），即可证明 \(Q(n)\)
对所有 \(n\) 成立。
\end{taplauthorsolutionentrybox}

\translatornote{原书只给出了上面的证明提示，并略去了具体的归纳步骤。下面
分别展开三条归纳原理。

\medskip\noindent
\textbf{（1）对深度归纳。}
题设假定：对于任意项 \(s\)，只要所有满足
\(\termdepth(r)<\termdepth(s)\) 的项 \(r\) 都有 \(P(r)\)，就能推出
\(P(s)\)。我们的目标是证明所有项 \(s\) 都满足 \(P(s)\)。

定义自然数上的命题
\[
  Q(n)\quad\Longleftrightarrow\quad
  \text{每个深度为 \(n\) 的项 \(s\) 都满足 \(P(s)\)。}
\]
依照原书的提示，我们对自然数 \(n\) 作完全归纳。任取 \(n\in\mathbb{N}\)，
假设 \(Q(m)\) 对所有满足 \(m<n\) 的自然数 \(m\) 都成立；需要证明
\(Q(n)\)。

若 \(n=0\)，则不存在深度为 \(0\) 的项，所以 \(Q(0)\) 空真。这是完全归纳
中的基础情形；它本身不包含任何项。

以下设 \(n\geq1\)。

任取一个满足 \(\termdepth(s)=n\) 的项 \(s\)。为了利用题设假定，需要证明：
每个满足 \(\termdepth(r)<\termdepth(s)\) 的项 \(r\) 都满足 \(P(r)\)。
于是任取满足
\[
  \termdepth(r)<\termdepth(s)=n
\]
的项 \(r\)，并令 \(m=\termdepth(r)\)。于是 \(m\) 是满足 \(m<n\) 的
自然数，归纳假设 \(Q(m)\) 因而适用，并给出 \(P(r)\)。

由于 \(r\) 是任取的，我们已经证明了上述条件中的每个项 \(r\) 都满足 \(P(r)\)。
而题设假定正是：这一条件成立时可以推出 \(P(s)\)。因此 \(P(s)\) 成立。
又因为 \(s\) 是任取的，所以所有深度为 \(n\) 的项都满足 \(P\)，即
\(Q(n)\) 成立。

特别地，当 \(n=1\) 时，不存在满足
\(\termdepth(r)<\termdepth(s)\) 的项 \(r\)，上述条件自动满足；题设假定
于是直接给出 \(P(s)\)。也就是说，\(n=0\) 是归纳原理的基础情形，而
\(n=1\) 是第一个包含实际项的层次。

完全归纳原理最终说明 \(Q(n)\) 对所有自然数 \(n\) 都成立。每个项都有某个
有限的自然数深度，因此 \(P(s)\) 对所有项 \(s\) 都成立。

\medskip\noindent
\textbf{（2）对大小归纳。}
题设假定：对于任意项 \(s\)，只要所有满足
\(\termsize(r)<\termsize(s)\) 的项 \(r\) 都有 \(P(r)\)，就能推出
\(P(s)\)。定义
\[
  R(n)\quad\Longleftrightarrow\quad
  \text{每个大小为 \(n\) 的项 \(s\) 都满足 \(P(s)\)。}
\]
对自然数 \(n\) 作完全归纳。任取 \(n\in\mathbb{N}\)，假设 \(R(m)\) 对
所有满足 \(m<n\) 的自然数 \(m\) 都成立，并证明 \(R(n)\)。若 \(n=0\)，
不存在大小为 \(0\) 的项，所以 \(R(0)\) 空真。以下设 \(n\geq1\)，并任取
大小为 \(n\) 的项 \(s\)。

若 \(\termsize(r)<\termsize(s)\)，令 \(m=\termsize(r)\)，则 \(m<n\)。
由归纳假设 \(R(m)\) 可得 \(P(r)\)。由于 \(r\) 是任取的，
每个满足 \(\termsize(r)<\termsize(s)\) 的项 \(r\) 都满足 \(P(r)\)。
根据题设假定，这就推出 \(P(s)\)，所以
\(R(n)\) 成立。

当 \(n=1\) 时，不存在大小更小的项，题设假定直接给出 \(P(s)\)。完全归纳
原理于是说明 \(P(s)\) 对所有项 \(s\) 都成立。这里同样是 \(n=0\) 构成
空的基础情形，\(n=1\) 才是第一个包含实际项的层次。

\medskip\noindent
\textbf{（3）结构归纳。}
题设假定：对于任意项 \(s\)，只要 \(s\) 的所有直接子项都满足 \(P\)，就能
推出 \(P(s)\)。定义
\[
  U(n)\quad\Longleftrightarrow\quad
  \text{每个深度不大于 \(n\) 的项都满足 \(P\)。}
\]
下面对自然数 \(n\) 作普通归纳。

基础情形是 \(n=0\)。不存在深度不大于 \(0\) 的项，所以 \(U(0)\) 空真。
这并不是在证明某个实际项满足 \(P\)，而是在为从 \(U(0)\) 推出 \(U(1)\)
准备归纳起点。

对于归纳步骤，假设 \(U(n)\) 成立，需要证明 \(U(n+1)\)。任取深度不大于
\(n+1\) 的项 \(s\)。若 \(\termdepth(s)\leq n\)，则由归纳假设
\(U(n)\) 直接得到 \(P(s)\)。否则 \(\termdepth(s)=n+1\)。根据深度的
定义，\(s\) 的每个直接子项的深度都不大于 \(n\)，所以归纳假设说明这些
直接子项都满足 \(P\)。再由题设假定推出 \(P(s)\)。

因此，每个深度不大于 \(n+1\) 的项都满足 \(P\)，即 \(U(n+1)\) 成立。
普通自然数归纳最终说明 \(P(s)\) 对所有项 \(s\) 都成立。}

\taplsolutionbacklink{thm:3.3.4}{返回定理 3.3.4}

\begin{taplauthorsolutionentrybox}
\phantomsection
\label{sol:author:3.5.5}
\textbf{3.5.5 解答}

设 \(P\) 是求值陈述之推导上的谓词。若对每个推导 \(\mathcal D\)，在已知
\(P(\mathcal C)\) 对 \(\mathcal D\) 的所有直接子推导
\(\mathcal C\) 成立时，都能证明 \(P(\mathcal D)\)，那么
\(P(\mathcal D)\) 对所有推导 \(\mathcal D\) 成立。

\taplsolutionbacklink{ex:3.5.5}{返回练习 3.5.5}
\end{taplauthorsolutionentrybox}

\begin{taplauthorsolutionentrybox}
\phantomsection
\label{sol:author:3.5.10}
\textbf{3.5.10 解答}

\[
\inferrule{t\trans t'}{t\transmulti t'}
\qquad
\inferrule{ }{t\transmulti t}
\qquad
\inferrule{t\transmulti t'\\t'\transmulti t''}
          {t\transmulti t''}
\]

\taplsolutionbacklink{ex:3.5.10}{返回练习 3.5.10}
\end{taplauthorsolutionentrybox}

\begin{taplauthorsolutionentrybox}
\phantomsection
\label{sol:author:3.5.13}
\textbf{3.5.13 解答}

\begin{enumerate}[label=\textup{(\arabic*)}]
  \item 定理 3.5.4 与定理 3.5.11 不再成立；定理 3.5.7、3.5.8
    和 3.5.12 仍然成立。
  \item 现在只有定理 3.5.4 不成立，其余各条仍然成立。值得注意的是：
    尽管加入这条规则后，单步求值变得不确定，多步求值的最终结果仍然是
    确定的，也就是说“条条大路通罗马”。严格证明并不困难，只是没有原来
    那么直接。关键在于单步求值关系具有下面这项局部汇合性质。
\end{enumerate}

\medskip
\textbf{引理 A.1（弱菱形性质）。}
若 \(r\trans s\)、\(r\trans t\) 且 \(s\neq t\)，则以下三种情形至少有
一种成立：
\[
  s\trans t,\qquad
  t\trans s,\qquad
  \text{存在 \(u\)，使 \(s\trans u\) 且 \(t\trans u\)。}
\]

\textit{证明。}
从求值规则可见，只有当
\(r=\taplif{r_1}{r_2}{r_3}\) 时才会出现这种情形。对
\(r\trans s\) 与 \(r\trans t\) 的两个推导同时作归纳，并对二者最后
使用的规则分类讨论。

例如，设 \(r\trans s\) 使用 \textsc{E-IfTrue}，
\(r\trans t\) 使用 \textsc{E-Funny2}。由规则形式可知
\[
  s=r_2,\qquad
  t=\taplif{\tapltrue}{r_2'}{r_3},
  \qquad r_2\trans r_2'
\]
取 \(u=r_2'\) 即可，因为 \(s\trans u\)，而
\(t\trans u\) 可由 \textsc{E-IfTrue} 得到。

若第一个推导使用 \textsc{E-IfFalse}、第二个使用
\textsc{E-Funny2}，则
\[
  s=r_3,\qquad
  t=\taplif{\taplfalse}{r_2'}{r_3},
\]
此时直接有 \(t\trans s\)。

再如，若两个推导最后都使用 \textsc{E-If}，则
\[
\begin{aligned}
s&=\taplif{r_1'}{r_2}{r_3},&
t&=\taplif{r_1''}{r_2}{r_3},\\
r_1&\trans r_1',&
r_1&\trans r_1''
\end{aligned}
\]
由归纳假设，\(r_1'\) 与 \(r_1''\) 或者能直接归约到彼此，或者能各自
归约到某个共同的 \(r_1'''\)。在前两种情形中，相应的两个条件式也能直接
归约到彼此；在后一种情形中，取
\[
  u=\taplif{r_1'''}{r_2}{r_3}
\]
即可。其他情形类似。\qed

\medskip
结果唯一性可由一次直接的“图追踪”得到。若
\(r\transmulti s\) 且 \(r\transmulti t\)，从两条序列的第一步开始，
反复应用引理 A.1：若一边一步到达另一边，就沿较长的一边继续；否则，把两个
后继拉到一个共同后继。如此逐层把两条路径拉拢起来。由于求值终止
（定理 3.5.12 仍然成立），这个过程最终得到两条路径的共同后继。

因此，若 \(r\) 既能求值到值 \(v\)，又能求值到值 \(w\)，则
\(v=w\)：二者都是范式，而两个范式要继续求值到同一个项，唯一的可能就是
它们本来相同。

\translatornote{原书附录 A 的 Lemma A.1 把这里的性质写成了更强的
菱形性质，但作者官方勘误指出该陈述为假；本解答采用勘误给出的弱菱形版本，
并相应调整了后面的图追踪论证。}

\taplsolutionbacklink{ex:3.5.13}{返回练习 3.5.13}
\end{taplauthorsolutionentrybox}

\begin{taplauthorsolutionentrybox}
\phantomsection
\label{sol:author:3.5.14}
\textbf{3.5.14 解答}

对 \(t\) 的结构作归纳。

\medskip\noindent
\textbf{情形：\(t\) 是值。}
每个值都处于范式，所以这一情形不可能发生。

\medskip\noindent
\textbf{情形：\(t=\taplsucc{t_1}\)。}
检查求值规则可知，只有 \textsc{E-Succ} 能用来导出
\(t\trans t'\) 与 \(t\trans t''\)：其他规则左侧的最外层构造子都不是
\(\mathsf{succ}\)。因此，必有两个子推导，其结论分别为
\(t_1\trans t_1'\) 与 \(t_1\trans t_1''\)。因为 \(t_1\) 是 \(t\)
的子项，归纳假设给出 \(t_1'=t_1''\)，于是
\(\taplsucc{t_1'}=\taplsucc{t_1''}\)。

\medskip\noindent
\textbf{情形：\(t=\taplpred{t_1}\)。}
可以把 \(t\) 归约为 \(t'\) 或 \(t''\) 的规则有三条：
\textsc{E-Pred}、\textsc{E-PredZero} 与 \textsc{E-PredSucc}。
不过，这三条规则彼此不重叠：若 \(t\) 匹配其中一条的左侧，就不可能再匹配
另外两条。例如，若它匹配 \textsc{E-Pred}，则 \(t_1\) 必定不是值，
特别不可能是 \(0\) 或 \(\mathsf{succ}\,v\)。所以两个推导最后必定使用
同一条规则。若该规则是 \textsc{E-Pred}，就像上一情形一样使用归纳假设；
若是另两条规则之一，结论立即得到。

\medskip\noindent
\textbf{其他情形。}
类似。

\taplsolutionbacklink{ex:3.5.14}{返回练习 3.5.14}
\end{taplauthorsolutionentrybox}

\begin{taplauthorsolutionentrybox}
\phantomsection
\label{sol:author:3.5.16}
\textbf{3.5.16 解答}

用元变量 \(t\) 表示加入 \(\taplwrong\) 后的新项集中的项（包括以
\(\taplwrong\) 为子短语的项），用 \(g\) 表示原来不含
\(\taplwrong\) 的“良好”项。记 \(\wtrans\) 为加入
\(\taplwrong\) 规则后的求值关系，\(\otrans\) 为原求值关系。两种处理方式
“彼此一致”可以精确表述为：一个原始项在旧语义中求值到卡住项，当且仅当
它在新语义中求值到 \(\taplwrong\)。

\medskip
\textbf{命题 A.2。}
对每个原始项 \(g\)，
\[
  \bigl(g\otransmulti g'\text{，其中 \(g'\) 卡住}\bigr)
  \quad\Longleftrightarrow\quad
  g\wtransmulti\taplwrong
\]
此外，对每个原始项 \(g\) 和原始值 \(v\)，
\[
  g\otransmulti v
  \quad\Longleftrightarrow\quad
  g\wtransmulti v
\]

证明分几步进行。首先，新增的归约规则不会破坏练习 3.5.14 中的确定性
性质。

\medskip
\textbf{引理 A.3。}
扩充后的求值关系是确定的。

这意味着：只要 \(g\) 在原语义中能单步归约到 \(g'\)，它在扩充语义中
也能归约到 \(g'\)，而且 \(g'\) 是扩充语义中 \(g\) 唯一能归约到的项。

其次证明，在原语义中已经卡住的项，在扩充语义中总会求值到
\(\taplwrong\)。

\medskip
\textbf{引理 A.4。}
若 \(g\) 卡住，则 \(g\wtransmulti\taplwrong\)。

\textit{证明。}
对 \(g\) 的结构作归纳。

\medskip\noindent
\textbf{情形：\(g=\tapltrue,\taplfalse\) 或 \(\taplzero\)。}
不可能，因为假定了 \(g\) 卡住。

\medskip\noindent
\textbf{情形：\(g=\taplif{g_1}{g_2}{g_3}\)。}
由于 \(g\) 卡住，\(g_1\) 必须处于范式，否则 \textsc{E-If} 可以应用。
\(g_1\) 显然不能是 \(\tapltrue\) 或 \(\taplfalse\)，否则
\textsc{E-IfTrue} 或 \textsc{E-IfFalse} 可以应用，\(g\) 就不会卡住。
考虑其余情形。

若 \(g_1=\taplif{g_{11}}{g_{12}}{g_{13}}\)，则 \(g_1\) 处于范式而又
不是值，所以它本身卡住。归纳假设给出
\(g_1\wtransmulti\taplwrong\)。由此可以构造
\[
  g\wtransmulti\taplif{\taplwrong}{g_2}{g_3},
\]
再使用 \textsc{E-If-Wrong} 得到
\(g\wtransmulti\taplwrong\)。

若 \(g_1=\taplsucc{g_{11}}\)，且 \(g_{11}\) 是数值，则 \(g_1\)
属于 \(\mathit{badbool}\)，规则 \textsc{E-If-Wrong} 立即给出
\(g\wtrans\taplwrong\)。否则，按定义 \(g_1\) 卡住，归纳假设给出
\(g_1\wtransmulti\taplwrong\)。从这个推导可构造
\[
  g\wtransmulti\taplif{\taplwrong}{g_2}{g_3},
\]
最后使用 \textsc{E-If-Wrong} 得到所需结果。其他子情形类似。

\medskip\noindent
\textbf{情形：\(g=\taplsucc{g_1}\)。}
因为 \(g\) 卡住，由值的定义可知，\(g_1\) 必须处于范式且不是数值。
于是有两种可能：\(g_1\) 是 \(\tapltrue\) 或 \(\taplfalse\)；或者
\(g_1\) 不是值，因而自身卡住。第一种情况下，
\textsc{E-Succ-Wrong} 立即给出
\(g\wtrans\taplwrong\)；第二种情况下，归纳假设给出
\(g_1\wtransmulti\taplwrong\)，随后按上面同样的方法完成推导。

\medskip\noindent
\textbf{其他情形。}
类似。\qed

引理 A.3 与 A.4 一起给出命题 A.2 从左到右的方向。反方向需要说明：
在扩充语义中“即将出错”的原始项，在原语义中必定卡住。

\medskip
\textbf{引理 A.5。}
若 \(g\wtrans t\)，而 \(t\) 含有 \(\taplwrong\) 作为子项，则
\(g\) 在原语义中卡住。

\textit{证明。}
对扩充语义中的求值推导作直接归纳即可。\qed

若 \(g\wtransmulti\taplwrong\)，考察序列中第一次出现
\(\taplwrong\) 的那一步。此前的每一步都是原语义中已有的求值步骤；设此前
到达的原始项为 \(h\)，则 \(g\otransmulti h\)，而引理 A.5 说明 \(h\)
在原语义中卡住。这就得到命题 A.2 从右到左的方向。关于原始值 \(v\) 的
第二个等价式也由同样的规则包含关系与确定性直接得到：求值序列一旦使用了
产生 \(\taplwrong\) 的规则，就不可能再以原始值结束。至此证明完毕。

\translatornote{作者官方勘误要求在命题 A.2 中补入关于原始值 \(v\) 的
第二个等价式，并修正引理 A.4 条件式子情形中的一段推导；这里已经一并落实。}

\taplsolutionbacklink{ex:3.5.16}{返回练习 3.5.16}
\end{taplauthorsolutionentrybox}

\begin{taplauthorsolutionentrybox}
\phantomsection
\label{sol:author:3.5.17}
\textbf{3.5.17 解答}

分别证明等价关系的两个方向。每个方向都先建立一个关于多步小步求值的技术
引理。

\medskip
\textbf{引理 A.6。}
若 \(t_1\transmulti t_1'\)，则
\[
  \taplif{t_1}{t_2}{t_3}
  \transmulti
  \taplif{t_1'}{t_2}{t_3}
\]
其他项构造子也有类似结论。

\textit{证明。}
简单归纳。\qed

\medskip
\textbf{命题 A.7。}
若 \(t\bigstep v\)，则 \(t\transmulti v\)。

\textit{证明。}
对 \(t\bigstep v\) 的推导作归纳，并按最后使用的规则分类。

\medskip\noindent
\textbf{情形 \rulename{B-Value}：}
\(t=v\)，结论立即成立。

\medskip\noindent
\textbf{情形 \rulename{B-IfTrue}：}
\[
  t=\taplif{t_1}{t_2}{t_3},
  \qquad t_1\bigstep\tapltrue,
  \qquad t_2\bigstep v
\]
由归纳假设，\(t_1\transmulti\tapltrue\)。再由引理 A.6，
\[
  \taplif{t_1}{t_2}{t_3}
  \transmulti
  \taplif{\tapltrue}{t_2}{t_3}
\]
使用 \textsc{E-IfTrue} 有
\(\taplif{\tapltrue}{t_2}{t_3}\trans t_2\)；另一次使用归纳假设得到
\(t_2\transmulti v\)。最后利用 \(\transmulti\) 的传递性即可。
其他情形类似。\qed

\medskip
\textbf{引理 A.8。}
若
\[
  \taplif{t_1}{t_2}{t_3}\transmulti v,
\]
则或者
\[
  t_1\transmulti\tapltrue
  \quad\text{且}\quad
  t_2\transmulti v,
\]
或者
\[
  t_1\transmulti\taplfalse
  \quad\text{且}\quad
  t_3\transmulti v
\]
而且，这些关于 \(t_1\) 与 \(t_2\)（或 \(t_3\)）的求值序列都严格短于
给定序列。其他项构造子也有类似结论。

\textit{证明。}
对给定求值序列的长度作归纳。条件式不是值，所以序列至少有一步。按第一步
使用的最后一条规则分类即可。若第一步使用 \textsc{E-If}，则
\[
  \taplif{t_1}{t_2}{t_3}
  \trans
  \taplif{t_1'}{t_2}{t_3}
  \transmulti v,
  \qquad t_1\trans t_1'
\]
由归纳假设，\(t_1'\) 最终归约到 \(\tapltrue\) 且 \(t_2\) 归约到
\(v\)，或者 \(t_1'\) 最终归约到 \(\taplfalse\) 且 \(t_3\) 归约到
\(v\)。把起始的一步 \(t_1\trans t_1'\) 接到相应序列之前，就得到结论。
所得序列显然比原序列短。

若第一步使用 \textsc{E-IfTrue}，则
\[
  \taplif{\tapltrue}{t_2}{t_3}\trans t_2\transmulti v,
\]
结论立即成立。\textsc{E-IfFalse} 情形同理，此时使用的是
\(t_3\transmulti v\)。\qed

\medskip
\textbf{命题 A.9。}
若 \(t\transmulti v\)，则 \(t\bigstep v\)。

\textit{证明。}
先证明辅助结论：若 \(t\trans t'\) 且 \(t'\bigstep v\)，则
\(t\bigstep v\)。对 \(t\trans t'\) 的推导作归纳，并按最后使用的
小步规则分类；每一情形都可以用相应的大步规则及归纳假设重建
\(t\bigstep v\)。

现在对 \(t\transmulti v\) 的步数作归纳。零步时 \(t=v\)，由
\textsc{B-Value} 得到结论。若序列至少有一步，把它写成
\[
  t\trans t'\transmulti v
\]
由归纳假设 \(t'\bigstep v\)，再应用刚才的辅助结论，得到
\(t\bigstep v\)。\qed

\translatornote{原书对命题 A.9 直接按小步序列长度归纳，但作者官方勘误
指出该证明在 \(\mathsf{succ}\) 情形中不成立，因为子项与整项的求值步数
可能相同。这里采用勘误建议的辅助引理修复证明。}

\taplsolutionbacklink{ex:3.5.17}{返回练习 3.5.17}
\end{taplauthorsolutionentrybox}
